\begin{titlepage}
    \begin{abstract}
        La gestione tradizionale di contest è un processo frammentato e inefficiente, caratterizzato dall'uso di strumenti generici non integrati, che generano criticità in termini di sicurezza, integrità dei dati e trasparenza.
        Il presente lavoro di tesi affronta questa problematica attraverso la progettazione, l'implementazione e la distribuzione di Swift Contest, una piattaforma software completa e multipiattaforma.

        Il sistema è stato sviluppato adottando un'architettura client-server moderna, con un frontend realizzato in Flutter e un backend basato su Supabase, che offre una suite completa di servizi gestiti, dall'autenticazione allo storage, fino all'esecuzione di logica serverless.
        La progettazione è stata guidata da principi di ingegneria del software quali la Separazione delle Responsabilità (SoC) e il Principio del Minimo Privilegio (PoLP), garantendo un accesso ai dati mediato da uno strato di API sicure.

        Swift Contest dimostra la fattibilità di costruire un sistema complesso, sicuro e scalabile con un ecosistema tecnologico accessibile, validando l'efficacia di un approccio a codebase singola e l'ottimizzazione dei costi operativi.
        Il prototipo finale si presenta come una soluzione integrata e robusta, pronta per una fase di validazione sul campo.
        Esso risolve efficacemente le inefficienze dei processi tradizionali e funge da solida base per futuri sviluppi nel campo delle piattaforme collaborative.
        \\[1cm]
    \end{abstract}
\end{titlepage}